\section*{Hoofdstuk \ref{ch:b2:atomic-orbitals} --- Atoomorbitalen: vormen, energieën en afmetingen}

\begin{solution}{exo:b2:atomic-orbitals:1}
3p: $n = 3$, $l = 1$; één \omterm{def:b2:atomic-orbitals:node}{radiale knoop} ($n - l - 1$) en één \omterm{def:b2:atomic-orbitals:node}{hoekknoop}. 4d: $n = 4$,
$l = 2$; één \omterm{def:b2:atomic-orbitals:node}{radiale knoop} en twee hoekknopen.
\end{solution}

\begin{solution}{exo:b2:atomic-orbitals:2}
$P_{2p} \propto r^4\mathrm e^{-r}$ ($r$ in $a_0$): $d\ln P/dr = 4/r - 1 = 0$ voor $r = 4$:
$4a_0 = \qty{212}{pm}$.
\end{solution}

\begin{solution}{exo:b2:atomic-orbitals:3}
$1 - \mathrm e^{-2}(1 + 2 + 2) = 1 - 5\mathrm e^{-2} = 0.323$.
\end{solution}

\begin{solution}{exo:b2:atomic-orbitals:4}
Vier lobben die tussen de $x$- en $y$-as wijzen, in het $xy$-vlak, met afwisselend
teken (positief waar $xy > 0$); de knoopvlakken zijn $xz$ en $yz$. Radiale knopen:
$n - l - 1 = 0$.
\end{solution}

\begin{solution}{exo:b2:atomic-orbitals:5}
Zuurstof, $(1\mathrm s)^2(2\mathrm s, 2\mathrm p)^6$: $\sigma = 5 \times 0.35 + 2 \times 0.85
= 3.45$, $Z^* = 4.55$; straal $4/4.55 = 0.88a_0$, \qty{47}{pm}.
\end{solution}

\begin{solution}{exo:b2:atomic-orbitals:6}
Ionisatie verwijdert de minst gebonden elektronen. De 4s-elektronen
(\qty{-14}{eV}) zijn veel minder sterk gebonden dan de 3d-elektronen (\qty{-59}{eV}): beide 4s-elektronen
vertrekken eerst, daarna één 3d-elektron, wat $[\ce{Ar}]\,3\mathrm d^5$ geeft.
\end{solution}

\begin{solution}{exo:b2:atomic-orbitals:7}
Lithium, 2s: $Z^* = 3 - 2 \times 0.85 = 1.30$, $I = 13.6 \times 1.30^2/4 =
\qty{5.75}{eV}$ (gemeten 5.39, 7~\% te hoog). Natrium, 3s: $Z^* = 11 - 8 \times 0.85 -
2 \times 1.00 = 2.20$, $I = 13.6 \times 2.20^2/9 = \qty{7.3}{eV}$ (gemeten 5.14): de
regels, een ruw model, overschatten de binding van het buitenste elektron, en steeds meer
naarmate men in de groep afdaalt.
\end{solution}

\begin{solution}{exo:b2:atomic-orbitals:8}
Natrium: $Z^* = 2.20$, straal $9/2.20 = 4.1a_0$ (\qty{216}{pm}). Chloor, 3p:
$\sigma = 6 \times 0.35 + 8 \times 0.85 + 2 \times 1.00 = 10.90$, $Z^* = 6.10$, straal
$9/6.10 = 1.5a_0$ (\qty{78}{pm}). Dezelfde schil, maar het buitenste elektron van chloor
voelt bijna driemaal de lading, en wordt naar binnen getrokken.
\end{solution}

\begin{solution}{exo:b2:atomic-orbitals:9}
Beide hebben het \omterm{def:b2:atomic-orbitals:radial-angular}{hoekdeel} $1/\sqrt{4\pi}$, dus de integraal herleidt tot
\[
  \int_0^\infty R_{1s}R_{2s}r^2\,dr \propto \int_0^\infty (2 - r)r^2\mathrm e^{-3r/2}\,dr =
  2 \times \frac{2!}{(3/2)^3} - \frac{3!}{(3/2)^4} = \frac{32}{27} - \frac{32}{27} = 0 .
\]
\end{solution}

\begin{solution}{exo:b2:atomic-orbitals:10}
$\int_0^\infty N^2r^2\mathrm e^{-r}r^2\,dr = N^2 \times 4! = 1$: $N = 1/\sqrt{24} =
1/(2\sqrt6)$, zoals in de tabel.
\end{solution}

\begin{solution}{exo:b2:atomic-orbitals:11}
$\langle r\rangle = 1.5a_0$, \omterm{def:b2:atomic-orbitals:radial-distribution}{meest waarschijnlijke straal} $a_0$. $P(r) = 4r^2\mathrm e^{-2r}$ stijgt
steil vanaf nul en daalt langzaam: de lange staart bij grote $r$ weegt zwaarder in het
gemiddelde dan in de positie van het maximum.
\end{solution}

\begin{solution}{exo:b2:atomic-orbitals:12}
$Y_{p_z} \propto z/r = \cos\theta$ wordt nul voor $\theta = \pi/2$: het $xy$-vlak.
$\psi$ heeft het teken van $z$, positief erboven, negatief eronder. Op het vlak is de
\omterm{def:b2:atomic-orbitals:wavefunction}{waarschijnlijkheidsdichtheid} nul.
\end{solution}

\begin{solution}{pb:b2:atomic-orbitals:1}
\textbf{1.} $\int|\psi|^2\,dV = \frac{4\pi}{\pi a_0^3}\int_0^\infty r^2\mathrm e^{-2r/a_0}\,dr =
\frac{4}{a_0^3} \times \frac{a_0^3}{4} = 1$.
\textbf{2.} $P(r) = 4\pi r^2|\psi|^2 = (4/a_0^3)\,r^2\mathrm e^{-2r/a_0}$.
\textbf{3.} $dP/dr = 0$: $2/r - 2/a_0 = 0$, $r = a_0$.
\textbf{4.} $\langle r\rangle = (4/a_0^3)\int_0^\infty r^3\mathrm e^{-2r/a_0}\,dr = (4/a_0^3)
\times 6(a_0/2)^4 = 1.5a_0$.
\textbf{5.} De verdeling is scheef, met een lange staart bij grote $r$.
\textbf{6.} $|\psi|^2$ is een dichtheid per volume-eenheid, het grootst bij de kern; $P(r)$
telt het volume van de schil mee, $4\pi r^2\,dr$, dat nul wordt bij $r = 0$.
\textbf{7.} \qty{52.9}{pm}.
\textbf{8.} Met $x = r/a_0$: $\int_{x_0}^\infty 4x^2\mathrm e^{-2x}\,dx$; tweemaal partieel geïntegreerd
$= \mathrm e^{-2x_0}(2x_0^2 + 2x_0 + 1)$.
\textbf{9.} $5\mathrm e^{-2} = 0.677$.
\textbf{10.} $13\mathrm e^{-4} = 0.238$.
\textbf{11.} $\mathrm e^{-2x}(1 + 2x + 2x^2) = 0.10$ voor $x \approx 2.66$: \qty{141}{pm}.
\textbf{12.} De dichtheid wordt op geen enkele afstand nul: geen rand. Maar 90~\% van de
waarschijnlijkheid ligt binnen \qty{141}{pm}: een praktische afmeting.
\textbf{13.} 2s: één \omterm{def:b2:atomic-orbitals:node}{radiale knoop} (bol met straal $2a_0$), geen \omterm{def:b2:atomic-orbitals:node}{hoekknoop}. 2p: één
\omterm{def:b2:atomic-orbitals:node}{hoekknoop} (een vlak), geen \omterm{def:b2:atomic-orbitals:node}{radiale knoop}.
\textbf{14.} Bij $r = 2a_0$.
\textbf{15.} $4a_0$ (oefening 2).
\textbf{16.} $d\ln P/dr = 2/r - 2/(2 - r) - 1 = 0$ geeft $r^2 - 6r + 4 = 0$, $r = 3 \pm \sqrt5$:
$0.76a_0$ (klein binnenste maximum) en $5.24a_0$ (hoofdmaximum).
\textbf{17.} 2s, via zijn binnenste maximum. In een atoom met meer elektronen dringt een 2s-elektron
door de 1s-schil heen, wordt het minder afgeschermd dan een 2p-elektron en ligt het lager in
energie.
\textbf{18.} Alle stralen worden gedeeld door $Z = 6$: 2p bij $0.67a_0$ (\qty{35}{pm}); maxima van 2s
bij $0.13a_0$ en $0.87a_0$.
\textbf{19.} C 3.25, N 3.90, O 4.55.
\textbf{20.} $\varepsilon = -13.6\,Z^{*2}/4$: C \qty{-35.9}{eV}, N \qty{-51.7}{eV}, O
\qty{-70.4}{eV}.
\textbf{21.} Zoals in het hoofdstuk: \qty{11.5}{eV} (gemeten 11.26).
\textbf{22.} N: atoom $5 \times (-13.6 \times 3.90^2/4) = \qty{-258.6}{eV}$; \ce{N+}, $Z^* =
4.25$, $4 \times (-13.6 \times 4.25^2/4) = \qty{-245.7}{eV}$; $I = \qty{12.9}{eV}$.
\textbf{23.} O: atoom $6 \times (-13.6 \times 4.55^2/4) = \qty{-422.3}{eV}$; \ce{O+}, $Z^* =
4.90$, $5 \times (-13.6 \times 4.90^2/4) = \qty{-408.2}{eV}$; $I = \qty{14.2}{eV}$. Model
11.5, 12.9, 14.2; gemeten 11.26, 14.53, 13.62.
\textbf{24.} De gemeten energie van zuurstof is lager dan die van stikstof: in
zuurstof moet het vierde 2p-elektron een orbitaal delen met een ander, en hun
afstoting maakt het gemakkelijker te verwijderen, terwijl de drie ongepaarde p-elektronen van stikstof
een stabiele, half gevulde subschil vormen. De \omterm{def:b2:atomic-orbitals:slater}{regels van Slater} behandelen alle elektronen van een groep
gelijk en weten niets van paarvorming.
\textbf{25.} Het 1s-elektron van waterstof bevindt zich verder dan $2a_0$ met waarschijnlijkheid
$\boldsymbol{13\mathrm e^{-4} \approx 0.238}$.
\end{solution}
